% Block 3 poster - Tasks 1 to 4.5
% Compile with pdflatex. Figures are read from the folder Plots/ :
%   task1_step.pdf, task2_pole_zero.pdf, task3_pi_window.pdf, task4_designs.pdf
\documentclass[8pt]{extarticle}
\usepackage[a3paper,portrait,margin=9mm]{geometry}
\usepackage[utf8]{inputenc}
\usepackage{amsmath,amssymb,bm}
\usepackage{graphicx}
\usepackage{booktabs}
\usepackage{xcolor}
\usepackage{enumitem}
\usepackage[most]{tcolorbox}
\usepackage{microtype}
\usepackage{multicol}
\setlength{\columnsep}{7mm}

\graphicspath{{Plots/}}
\definecolor{ntnu}{HTML}{00509E}
\definecolor{soft}{HTML}{EAF1F8}
\pagestyle{empty}
\setlength{\parindent}{0pt}
\setlength{\parskip}{2.5pt}
\setlist{leftmargin=1.2em,itemsep=1pt,topsep=2pt,parsep=0pt}

% numbered section bar, same look as the Block 2 poster
\newcounter{psec}
\newcommand{\psection}[1]{%
  \refstepcounter{psec}\vspace{5pt}%
  \colorbox{ntnu}{\parbox{\dimexpr\linewidth-2\fboxsep}{%
    \color{white}\sffamily\bfseries\large\thepsec\quad #1}}\par\vspace{2pt}}

% framed key result
\newtcolorbox{keybox}{colback=soft,colframe=ntnu,boxrule=0.5pt,arc=0pt,
  left=4pt,right=4pt,top=3pt,bottom=3pt,before skip=5pt,after skip=5pt}

\newcommand{\PO}{\mathrm{P.O.}}
\newcommand{\zwn}{\zeta\omega_n}
\newcommand{\ub}{\bar u}
\newcommand{\ib}{\bar\imath_L}
\newcommand{\vb}{\bar v_C}

\begin{document}

% ---------------------------------------------------------------- title


\begin{multicols}{2}\raggedcolumns
% ================================================================ TASK 1
\psection{Task 1: overshoot and peak time fix the poles}

Template closed loop, with $\zeta$ and $\omega_n$ free:
\[
G(s)=\frac{\omega_n^2}{s^2+2\zeta\omega_n s+\omega_n^2},
\qquad s_{1,2}=\alpha\pm j\beta=-\zwn\pm j\,\omega_n\sqrt{1-\zeta^2}.
\]
The poles equal to the eigenvalues $\alpha\pm j\beta$ from Block~2.
The slope of the step response is proportional to $e^{\alpha t}\sin\beta t$,
so the first peak is at $\beta t=\pi$.

\begin{keybox}
\[
T_p=\frac{\pi}{\beta}=\frac{\pi}{\omega_n\sqrt{1-\zeta^2}},
\qquad
\frac{\PO}{100}=e^{-\zwn T_p}=\exp\!\Big(\frac{-\pi\zeta}{\sqrt{1-\zeta^2}}\Big)
\]
\end{keybox}

\[
L=\ln\frac{\PO}{100}:\quad
\zeta=\frac{|L|}{\sqrt{\pi^2+L^2}},\qquad
\omega_n=\frac{\pi}{T_p\sqrt{1-\zeta^2}},\qquad
\zwn=\frac{|L|}{T_p}.
\]
$\PO$ sets $\zeta$ alone (the angle of the poles), $T_p$ then sets $\beta$
(their height), so any pair of specifications can be met. For 5\,\%, 1\,ms:
$\zeta=0.690$, $\omega_n=4341$\,rad/s, poles $-2996\pm j3142$.

\begin{minipage}[c]{0.42\linewidth}
\includegraphics[width=\linewidth]{task1_step.pdf}
\end{minipage}\hfill
\begin{minipage}[c]{0.50\linewidth}\centering
\small
\begin{tabular}{@{}rrrr@{}}
\toprule
\multicolumn{2}{c}{Specified} & \multicolumn{2}{c}{\texttt{stepinfo}}\\
\cmidrule(r){1-2}\cmidrule(l){3-4}
P.O. & $T_p$ & P.O. & $T_p$\\
\midrule
5\,\%  & 1\,ms & 5.00\,\%  & 0.999\,ms\\
5\,\%  & 3\,ms & 5.00\,\%  & 3.0\,ms\\
12\,\% & 1\,ms & 12.00\,\% & 0.999\,ms\\
18\,\% & 1\,ms & 18.00\,\% & 0.994\,ms\\
\bottomrule
\end{tabular}
\end{minipage}


% ================================================================ TASK 2
\psection{Task 2: when an extra pole or zero can be ignored}

\[
H(s)=\frac{(\gamma/a)(s+a)\,\omega_n^2}{(s^2+2\zeta\omega_n s+\omega_n^2)(s+\gamma)}
\]
The factor $\gamma/a$ gives $H(0)=1$ for any $\gamma$ and $a$, so only the
transient changes. Design fixed at 5\,\%, 1\,ms; $\gamma$ and $a$ in units of
$10\,\zwn$, one moved at a time.

\begin{minipage}[c]{0.42\linewidth}
\includegraphics[width=\linewidth]{task2_pole_zero.pdf}
\end{minipage}\hfill
\begin{minipage}[c]{0.50\linewidth}\centering
\small
\begin{tabular}{@{}rrrr@{}}
\toprule
$\dfrac{\gamma}{10\zwn}$ & $\dfrac{a}{10\zwn}$ & P.O. & $T_p$\\
\midrule
1   & 10  & 4.94\,\%  & 1.03\,ms\\
10  & 1   & 5.06\,\%  & 0.97\,ms\\
0.1 & 10  & 0\,\%     & no peak\\
10  & 0.1 & 23.4\,\%  & 0.51\,ms\\
0.1 & 0.1 & 5.00\,\%  & 1.00\,ms\\
\bottomrule
\end{tabular}
\end{minipage}

\begin{itemize}
\item At the $10\,\zwn$ boundary the error is 0.06 percentage
  points in overshoot and 3\,\% in peak time.
\item A close zero gives $y_H=y_G+\tfrac1a\,\dot y_G$: faster, more
  overshoot. A close pole is a low-pass filter: slower, no overshoot.
\end{itemize}


% ================================================================ TASK 3
\psection{Task 3: a PI on a first-order plant leaves a narrow window}

Plant $G=\dfrac{k}{s+\alpha}$ with $k=-5$, $\alpha=20$; controller
$G_c=\kappa_P\dfrac{s+\gamma}{s}$. With $L=N/D$, unity feedback gives
$T=N/(D+N)$:
\[
T(s)=\frac{k\kappa_P\,(s+\gamma)}{s^2+(\alpha+k\kappa_P)\,s+k\kappa_P\gamma}
\]
Matching the denominator to $s^2+2\zwn s+\omega_n^2$:

\begin{keybox}
\[
\kappa_P=\frac{2\zwn-\alpha}{k},\qquad
\gamma=\frac{\omega_n^2}{2\zwn-\alpha}
\]
\end{keybox}

Two gains, two specifications: the zero is forced. It stays out of the way
only if $|\gamma|\ge 10\,\zwn$, which for $\PO=5\,\%$ gives
\[
13.1\le\omega_n\le16.2\ \text{rad/s}
\quad\Longleftrightarrow\quad
0.27\ \text{s}\le T_p\le0.33\ \text{s}.
\]
The window is centred on $\omega_n=14.49$, where $\kappa_P=0$: a pure
integrator, no zero. Below it the zero is in the right half-plane.

\begin{minipage}[c]{0.44\linewidth}
\includegraphics[width=\linewidth]{task3_pi_window.pdf}
\end{minipage}\hfill
\begin{minipage}[c]{0.47\linewidth}\centering
\small
\begin{tabular}{@{}rrrr@{}}
\toprule
$\omega_n$ & $\dfrac{\gamma}{10\zwn}$ & P.O. & $T_p$ [ms]\\
\midrule
15 & $3.1$  & 5.00\,\% & $289\to285$\\
14 & $-3.0$ & 5.00\,\% & $310\to315$\\
30 & $0.2$  & 7.72\,\% & $145\to109$\\
\bottomrule
\end{tabular}

\smallskip
$T_p$: designed $\to$ measured.
\end{minipage}

Inside the window the design holds. At $\omega_n=30$ the zero is close and
the overshoot is 7.7\,\% instead of 5\,\%.
% ================================================================ TASK 4.1
\psection{Task 4.1: transfer function of the converter}

Linearized model from Block~2 at $\vb=48$\,V ($\ib=13.82$\,A, $\ub=0.2365$),
state $\tilde x=[\tilde\imath_L,\ \tilde v_C]^{\!\top}$, input $\tilde u$:
\[
A=\begin{bmatrix}-\dfrac{R}{L} & -\dfrac{1-\ub}{L}\\[8pt]
\dfrac{1-\ub}{C} & -\dfrac{G}{C}\end{bmatrix},
\qquad
B=\begin{bmatrix}\dfrac{\vb}{L}\\[8pt]-\dfrac{\ib}{C}\end{bmatrix},
\qquad
C_y=\begin{bmatrix}1&0\end{bmatrix}.
\]
The inner loop controls the inductor current, so $y=\tilde\imath_L$.
Transforming $\dot{\tilde x}=A\tilde x+B\tilde u$ with zero initial conditions:
\[
(sI-A)X=BU\ \Rightarrow\ X=(sI-A)^{-1}BU,
\qquad
G_{iL}(s)=\frac{Y}{U}=C_y\,(sI-A)^{-1}B .
\]
With $(sI-A)^{-1}=\operatorname{adj}(sI-A)/\det(sI-A)$, and $C_y$ selecting the
first row $[\,s-a_{22},\ a_{12}\,]$ of the adjugate:

\begin{keybox}
\[
G_{iL}(s)=\frac{b_1 s+a_{12}b_2-a_{22}b_1}
{s^2-(a_{11}+a_{22})\,s+a_{11}a_{22}-a_{12}a_{21}}
\]
\end{keybox}

Poles. The denominator is $\det(sI-A)$, so the poles are the
eigenvalues of Block~2.
Zero. The numerator depends on $B$ and on the chosen output:
\[
z=a_{22}-a_{12}\frac{b_2}{b_1}
\quad\Longrightarrow\quad
\boxed{\,z=-\frac{G}{C}-\frac{(1-\ub)\,\ib}{C\,\vb}\,}
\]
$L$ cancels: the zero depends on $C$, $G$ and the operating point only, and
is in the left half-plane at any operating point.

With $L=0.5$\,mH, $C=150\,\mu$F, $R=98$\,m$\Omega$, $G=2.71$\,mS:
\[
G_{iL}(s)=\frac{9.60\cdot10^{4}\,s+1.424\cdot10^{8}}{s^2+214.1\,s+7.776\cdot10^{6}}
\]
Poles $-107\pm j2786$\,s$^{-1}$ $(=\operatorname{eig}A)$ and zero
$-18-1465=-1483$\,s$^{-1}$; MATLAB \texttt{tf(ss(A,B,C\_y,0))} gives the same.

% ================================================================ TASK 4.2
\psection{Task 4.2: can the zero be ignored?}

Against the plant poles ($\zwn=107$):
\[
|z|=1483\ >\ 10\,\zwn=1070 \qquad\text{passes, by a factor 1.4}
\]
\begin{itemize}
\item Replace $(s+a)$ by $a$ to keep the final value:
\end{itemize}
\[
G_{iL,\mathrm{simp}}(s)=\frac{1.424\cdot10^{8}}{s^2+214.1\,s+7.776\cdot10^{6}},
\qquad
G_{iL,\mathrm{simp}}(0)=G_{iL}(0)=18.31
\]

Against the closed-loop poles:
\begin{itemize}
\item In $T=N/(D+N)$ feedback moves the poles, not the zero: it stays at $-1483$.
\item The specification formulas describe the closed loop, so compare with the
  poles we place: $\zwn=3.0/T_p$ for $\PO=5\,\%$.
\end{itemize}

\begin{keybox}
\[
|z|\ge10\,\zwn=\frac{30}{T_p}
\quad\Longrightarrow\quad
T_p\ \ge\ \frac{30}{1483}\approx 20\ \text{ms}
\]
\end{keybox}

\begin{itemize}
\item The plant rings with a period of 2.26\,ms, so 20\,ms is too slow.
\item At 1\,ms, $10\,\zwn=29\,957$: the zero cannot be neglected.
\end{itemize}

Redesign for 1\,ms ($|z|$ must be 20 times larger):
\begin{itemize}
\item Increase $G$: needs about 2\,S, 5\,kW of loss in a 500\,W
  converter. That is not plausible.
\item Reduce $C$: $z\propto1/C$, operating point unchanged, so
  $C\approx7.4\,\mu$F.
\item Cost: 20 times the output ripple; $\omega_n$ up by $\sqrt{20}$
  (444\,Hz to 2\,kHz, a decade below the switching frequency); $\zeta$ from
  0.038 to 0.022.
\end{itemize}

\psection{Task 4.3: closed loop with the leaky PI}

Controller $G_c=\kappa_P\dfrac{s+\gamma}{s+p}$, with $\gamma=\kappa_I/\kappa_P$
($p=0$ is the ordinary PI). Plant
$G_{iL}=\dfrac{b_1(s+a)}{s^2+a_1s+a_0}$, with $K=b_1a$ and
\[
a_1=-(a_{11}+a_{22})=214.1,\quad
a_0=a_{11}a_{22}-a_{12}a_{21}=7.776\cdot10^{6},\quad
b_1=9.60\cdot10^{4},\quad a=1483 .
\]
From $T=N/(D+N)$:

\begin{keybox}
\small
\[
T_{\mathrm{simp}}=\frac{K\kappa_P(s+\gamma)}
{s^3+(a_1+p)s^2+(a_0+a_1p+K\kappa_P)s+a_0p+K\kappa_P\gamma}
\]
\[
T_{\mathrm{full}}=\frac{b_1\kappa_P(s+a)(s+\gamma)}
{s^3+(a_1+p+b_1\kappa_P)s^2+\big(a_0+a_1p+b_1\kappa_P(a+\gamma)\big)s+a_0p+b_1\kappa_Pa\gamma}
\]
\end{keybox}

\begin{itemize}
\item Both are third order, with a controller zero at $-\gamma$.
\item The $s^2$ coefficient sets the total damping.
\item Simplified: $a_1+p$. No $\kappa_P$, so damping must come from the leak.
\item Full: $a_1+p+b_1\kappa_P$. The plant zero lets $\kappa_P$ damp the loop.
\end{itemize}

\psection{Task 4.4: design and test}

Attempt 1: pole placement on the simplified model for 5\,\%, 1\,ms.
Poles at $-2996\pm j3142$ and $-r$, with $r=10\,\zwn=29\,957$:
\[
p=2\zwn+r-a_1,\quad
\kappa_P=\frac{\omega_n^2+2\zwn r-a_0-a_1p}{K},\quad
\gamma=\frac{r\omega_n^2-a_0p}{K\kappa_P}
\]
Result: $\kappa_P=1.284$, $\kappa_I=2013$, $\gamma=1567$, $p=35\,735$.

Design 2: pure PI ($p=0$) tuned on the full model for 5\,\%, 0.4\,ms.
Result: $\kappa_P=0.0628$, $\kappa_I=122$.

\begin{minipage}[c]{0.46\linewidth}
\includegraphics[width=\linewidth]{task4_designs.pdf}
\end{minipage}\hfill
\begin{minipage}[c]{0.48\linewidth}\centering
\small
\begin{tabular}{@{}llrrr@{}}
\toprule
Design & Plant & P.O. & $T_p$ & $y(\infty)$\\
\midrule
Attempt 1 & simp. & 73.1\,\% & 0.40\,ms & 0.508\\
Attempt 1 & full  & 66.0\,\% & 0.28\,ms & 0.508\\
Design 2  & full  & 5.00\,\% & 0.40\,ms & 1.000\\
Block 2   & full  & 1.44\,\% & 0.32\,ms & 1.000\\
\bottomrule
\end{tabular}

\smallskip
Block 2: $K_{p,i}=0.1$, $K_{i,i}=100$.
\end{minipage}

\begin{itemize}
\item Attempt 1 fails: the controller zero at 1567 is far inside
  $10\,\zwn$, and the leak halves the final value.
\item Design 2 works: 5\,\% overshoot, no leak, no steady-state error.
\item 0.4\,ms, not 1\,ms: with the plant zero at $-1483$, well-damped
  designs peak at 0.3--0.4\,ms.
\end{itemize}

\psection{Task 4.5: steady-state error of the leak}

For a unit step, $e_{ss}=1-T(0)$. Both models give
\[
e_{ss}=\frac{a_0\,p}{a_0\,p+b_1\kappa_P\,a\,\gamma}
\]
\begin{center}\small
\begin{tabular}{@{}lrrr@{}}
\toprule
 & $p$ [s$^{-1}$] & $e_{ss}$ & P.O. (full plant)\\
\midrule
Attempt 1 & 35\,735 & 49.2\,\% & 66.0\,\%\\
Design 2 & 0 & 0 & 5.00\,\%\\
Design 2 with a leak & 100 & 4.3\,\% & 8.6\,\%\\
Design 2 with a leak & 1\,000 & 31\,\% & 38\,\%\\
\bottomrule
\end{tabular}
\end{center}
We are not satisfied with the compromise. The outer voltage loop relies on
$i_L$ following its reference, and a 49\,\% error defeats that. The leak was
only needed because the simplified model has no other source of damping; on
the full model $p=0$ works and $e_{ss}=0$.

\end{multicols}
\end{document}
